% !TEX root = ../main.tex
\chapter{Glossary}\label{ch:glossary}
Each entry says what the word means in this app, in a sentence or three, and where you meet
it. \Cref{ch:core-ideas} explains the main ones at more length.

\begin{glossarylist}
\glossitem{recording}{Recording} One experiment's data file: every electrode's voltage over
the whole time it was measured, with its sampling rate and, where known, its unit. You pick
one at the top left of Explore's Corpus page; Settings › Datasets imports new ones.

\glossitem{channel}{Channel} One electrode's trace inside a recording. Each row of the
coverage map in Explore is a channel; the Signal page shows one channel at a time.

\glossitem{span}{Span} A stretch of time on one channel, given by a start and an end. You
choose one in Explore and send it to Analyse, where it is the source of a chain.

\glossitem{dataset}{Dataset} A recording as the app lists and names it, with its channels,
sampling rate and unit. Settings › Datasets lists them. \unverified{In the app the two words
are used for the same thing.}

\glossitem{block}{Block} One analysis step with one job, one type in and one type out, and a
few settings. You add blocks in Analyse through the \ui{Insert a stage} window; each has a
card in \cref{ch:blocks}.

\glossitem{chain}{Chain} Blocks in a row, each feeding the next, starting from a span. Built
and run on Analyse's Chain page.

\glossitem{step}{Step} One position in a chain, holding one block and its settings. Analyse
numbers them \textsf{01}, \textsf{02} and so on; also called a stage.

\glossitem{template}{Template} A chain saved without its recording and span, so the same
steps can be applied elsewhere. Saved with \ui{Save template} in Analyse; listed in Library ›
Templates; applied to many channels in Discovery.

\glossitem{run}{Run} One execution of a chain over one span (or one channel). Every run is
numbered, keeps its recipe, and writes its span-set output as detections. Analyse's
\ui{History} lists past runs.

\glossitem{job}{Job} Work the app does in the background while you carry on: a run, a
Discovery run across channels, a model training, a cluster script. The \ui{Jobs} button at
the foot of the rail counts the live ones and opens the Jobs page.

\glossitem{signal}{Signal} The type of a trace: one value per moment plus the sampling rate.
Every chain starts from one; cleaning blocks take a Signal and give one back.

\glossitem{span-set}{Span set} The type of a list of stretches of time, each optionally
labelled and scored. Detectors produce one; a chain that ends in one writes its spans as
detections.

\glossitem{window}{Window} A fixed-length slice of a signal, such as ten minutes. Windows
are the unit that clustering and classifiers work on, and the unit of the project's older
hand labels.

\glossitem{window-set}{Window set} The type of a signal cut into windows, often with a table
of measurements per window. Library › Window sets keeps saved ones for reuse.

\glossitem{scores}{Scores} The type of a running value, one number per moment of the signal,
lined up with it. A threshold turns Scores into a span set.

\glossitem{encoding}{Encoding} The type of a signal redrawn in another form: a picture per
window, or a string of letters. Symbolic and image encoders produce one.

\glossitem{grouping}{Grouping} The type of a label per window, saying which group or class
it belongs to. Made by clustering or from human labels; read by the classifier block.

\glossitem{model}{Model} The type of a trained classifier, carried as a pointer to its saved
file. Trained and compared on the Models pages.

\glossitem{detection}{Detection} A span a machine proposed: one mark written by one run.
Kept in its own table, apart from annotations, and judged in Review.

\glossitem{annotation}{Annotation} A span a person marked, with that person's verdict. Explore's
\ui{Take span for Review} writes one; the project's older hand labels are annotations too.

\glossitem{verdict}{Verdict} A person's judgement of one span, one of exactly five:
\ui{seed}, \ui{interesting}, \ui{not interesting}, \ui{artifact}, \ui{unsure}. Given in
Review with the keys \texttt{S}, \texttt{I}, \texttt{N}, \texttt{A}, \texttt{U}.

\glossitem{candidate}{Candidate} A machine proposal waiting for a person's verdict: a
detection, or a match from a seed search. Review shows one candidate at a time.

\glossitem{review-queue}{Review queue} A list of candidates to judge, with its own name and
progress. Made where the candidates come from: Analyse's \ui{Pass N to Review}, Discovery,
Explore's \ui{Take span for Review}. The queue chip at the top of Review switches between
them.

\glossitem{family}{Family} A group of Library members that look alike, made by grouping
them by shape. Library › Atlas shows families as cards; Library › Family shows one.

\glossitem{motif}{Motif} A shape that recurs in the recordings. The Library is the catalogue
of motifs found so far, each with every place it occurs.

\glossitem{exemplar}{Exemplar} The member chosen to stand for a shape, and the shape a seed
search looks for. The \ui{seed} verdict in Review makes a new one in the Library.

\glossitem{medoid}{Medoid} The member of a family closest, on average, to all the others:
its most typical member, picked by the computer. Shown beside the exemplar on Library ›
Family; usable as a seed.

\glossitem{seed}{Seed} Two related uses. As a verdict, \ui{seed} marks a span worth keeping
as an example. As a search input, a seed is the shape a seed search looks for: a Library
exemplar, a family's medoid, or a span taken in Explore.

\glossitem{seed-search}{Seed search} A search for spans that look like one seed, across the
channels and hours you choose, with a null drawn alongside. Run from Discovery › Seed
search; its matches go to Review.

\glossitem{null}{Null} What a method finds by chance alone, measured by running it on
surrogates. A result counts only if the real signal beats its null. \Cref{app:nulls} has the
detail.

\glossitem{surrogate}{Surrogate} A fake copy of a signal that keeps its general character
but scrambles any real pattern, for example by randomising the phases of its frequencies.
Made by \blockref{preprocessing.surrogate}; never stored as a recording, so nothing found on
one can enter the Library.

\glossitem{held-out}{Held out} Kept unseen for the final test. The recording
\code{M4_aug_concat_fs1.mat} is held out: every page and the runner refuse it, and the header's
\ui{M4 held out} chip says so. Models also hold out part of their data for testing.

\glossitem{blocked-split}{Blocked split} A way of dividing windows into training and test
sets by time, in whole blocks with a gap between them, rather than at random. Neighbouring
windows look alike, so a random split would let the test peek at the training data. Used
when training models (\cref{ch:models}).

\glossitem{matrix-profile}{Matrix profile} A running score that says, for each moment, how
far the shape starting there is from its nearest look-alike anywhere else in the span. Low
values mark shapes that repeat; high values mark one-offs. Computed by
\blockref{detection.matrix_profile}.

\glossitem{symbolisation}{Symbolisation (SAX)} Turning a signal into a string of letters,
one per short piece, by its level or its slope, so that patterns can be searched for like
words in text. SAX (symbolic aggregate approximation) is the family of methods; the app
has several variants, such as \blockref{detection.sax_dsax}.

\glossitem{spike}{Spike} A brief rise of the signal clearly above its local surroundings.
\blockref{detection.spike_v1} finds them by one fixed rule, and can find brief dips the same
way.

\glossitem{spike-train}{Spike train} A run of spikes close together in time, read as one
event. Also one of Review's shape classes (key \texttt{1}). \blockref{interrogation.trains}
groups spikes into trains.

\glossitem{inter-spike-interval}{Inter-spike interval} The time from one spike to the next.
\blockref{interrogation.intervals} measures them from the events a detector such as
\blockref{detection.spike_v1} hands on.

\glossitem{adjudication}{Adjudication} A person's verdict on a detection, stored in its own
table, apart from both detections and annotations. Review's \ui{writes adjudications} chip
names where your verdicts go.

\glossitem{artifact}{Artifact} A mark caused by the recording set-up, not the fungus: a loose
contact, or two channels picking up the same interference. One of the five verdicts; kept
out of training.

\glossitem{tag}{Tag} A word from the project's controlled vocabulary attached to a span,
such as \ui{sharkfin}. Tags describe; they are never a verdict. Settings › Vocabulary holds
the list.

\glossitem{class}{Class} In Review, an optional shape category given beside a verdict, with
the keys \texttt{1}--\texttt{4} (\ui{spike-train}, \ui{burst}, \ui{slow-drift},
\ui{plateau}) and \texttt{9} (\ui{electrode artifact}).

\glossitem{member}{Member} One occurrence of a Library shape: a span in some recording, on
some channel. A family's members are listed on Library › Family.

\glossitem{edge}{Edge} A measured link between two Library members: how far apart their
shapes are, by which distance, and, across channels, the lag between them.

\glossitem{recipe}{Recipe} The exact description of a run: the source, the blocks in order
and every setting. Its fingerprint (the recipe hash) identifies every result the run made.

\glossitem{side-input}{Side input} A second input a block needs besides the previous step's
output, such as the exemplar a seeded search looks for. Bound when you build the chain.

\glossitem{coverage-map}{Coverage map} Explore's grid of channels against time bins, shaded
by how many annotations, detections or disagreements fall in each bin.

\glossitem{sampling-rate}{Sampling rate} How many values a recording holds per second, in
hertz (Hz): \code{M2_aug_concat_fs1.mat}, for example, holds one value a second (1~Hz).
Explore shows each recording's rate in the recordings menu.

\glossitem{cluster-hpc}{Cluster (computing)} The university's shared computers, used for runs
too long for a desktop. The app writes a SLURM script for the cluster's queue and imports
the results; see \cref{ch:jobs}. Not to be confused with clustering, the grouping of similar
windows.

\glossitem{demo-data}{Demo data} Example data a page shows where its live version is not yet
wired. A dashed \ui{demo data} chip in the page header marks it.

\glossitem{change-point}{Change point} A moment where a signal's behaviour changes and stays
changed, such as a jump in its level or in how much it varies.
\blockref{detection.rupture} finds them; \blockref{detection.fluss} finds where a recording
moves from one regime to another.

\glossitem{robust-noise}{Noise width (robust)} An estimate of how much a signal jitters from
moment to moment that a few large events barely move. Several detectors set their thresholds
as a multiple of it, so the same setting adapts to quiet and noisy channels; each block's
entry in \cref{app:blocks} says which estimate it uses.

\glossitem{image-encoding}{Image encoding} A way of turning a window of signal into a picture,
by comparing every moment in it with every other, so that image methods such as a neural
network can read it. Gramian angular fields and recurrence plots are the two kinds the app
makes (\blockref{catalogue.gramian_gasf}, \blockref{catalogue.gramian_recurrence}).

\glossitem{sample-entropy}{Sample entropy} A number for how unpredictable a stretch of signal
is: low when its patterns repeat, high when they do not. \blockref{catalogue.multiscale_entropy}
measures it at several time scales.

\glossitem{checkpoint}{Checkpoint} A trained neural network saved to a file, which a block
loads to score new windows. \blockref{catalogue.cnn_score} names its checkpoint in its
settings; its entry says which ones are fit to use.
\end{glossarylist}
